\documentclass[11pt, letterpaper]{article}

\usepackage[margin=1in, headheight=22pt, headsep=12pt]{geometry}
\usepackage{titlesec}
\usepackage{enumitem}
\usepackage{xcolor}
\usepackage{hyperref}
\usepackage{fancyhdr}
\usepackage{microtype}
\usepackage{parskip}
\usepackage{booktabs}
\usepackage{sectsty}
\usepackage{amssymb}

% --- Color Definitions ---
\definecolor{accent}{HTML}{2C3E6B}
\definecolor{subaccent}{HTML}{5B7BA5}
\definecolor{rulecolor}{HTML}{CCCCCC}
\definecolor{darktext}{HTML}{1A1A1A}
\definecolor{greenzone}{HTML}{2E7D32}
\definecolor{yellowzone}{HTML}{F9A825}
\definecolor{redzone}{HTML}{C62828}
\definecolor{grayzone}{HTML}{616161}
\definecolor{lightbg}{HTML}{F5F5F5}

% --- Hyperlink Setup ---
\hypersetup{
    colorlinks=true,
    linkcolor=accent,
    urlcolor=subaccent,
    pdftitle={Week 2: Big Bands and the Swing Era - Quiz Study Guide},
    pdfauthor={MUS 262},
}

% --- Section Formatting ---
\titleformat{\section}
    {\Large\bfseries\color{accent}}
    {\thesection}{1em}{}
    [\vspace{-0.5em}\textcolor{rulecolor}{\rule{\textwidth}{0.4pt}}]

\titleformat{\subsection}
    {\large\bfseries\color{subaccent}}
    {\thesubsection}{1em}{}

\titleformat{\subsubsection}
    {\normalsize\bfseries\color{accent}}
    {\thesubsubsection}{1em}{}

\titlespacing*{\section}{0pt}{1.5em}{0.75em}
\titlespacing*{\subsection}{0pt}{1em}{0.5em}
\titlespacing*{\subsubsection}{0pt}{0.75em}{0.4em}

% --- Header/Footer ---
\pagestyle{fancy}
\fancyhf{}
\renewcommand{\headrulewidth}{0.4pt}
\renewcommand{\headrule}{\hbox to\headwidth{\color{rulecolor}\leaders\hrule height \headrulewidth\hfill}}
\fancyhead[L]{\small\color{gray}MUS 262 --- Week 2: Big Bands \& Swing}
\fancyhead[R]{\small\color{gray}Quiz Study Guide}
\fancyfoot[C]{\small\color{gray}\thepage}

% --- List Formatting ---
\setlist[itemize]{leftmargin=1.5em, itemsep=3pt, parsep=0pt, topsep=4pt}
\setlist[itemize,2]{leftmargin=1.5em, itemsep=2pt}
\setlist[enumerate]{leftmargin=1.5em, itemsep=3pt, parsep=0pt, topsep=4pt}

% --- Custom Commands ---
\newcommand{\listenbox}[1]{\vspace{0.5em}\noindent\fcolorbox{rulecolor}{lightbg}{\parbox{\dimexpr\textwidth-2\fboxsep-2\fboxrule}{\small\textbf{Listen For:} #1}}\vspace{0.5em}}

% --- Body Text ---
\color{darktext}

\begin{document}

% --- Title Block ---
\thispagestyle{empty}
\begin{center}
    \vspace*{0.5cm}
    {\Huge\bfseries\color{accent} Week 2: Big Bands \&\\[0.3em] the Swing Era}\\[0.5em]
    {\Large\color{subaccent} Quiz Study Guide}\\[1em]
    \textcolor{rulecolor}{\rule{0.6\textwidth}{0.8pt}}\\[0.8em]
    {\color{gray}\small MUS/AAS 262 --- Jazz History: Many Sounds, Many Voices\\[0.3em]
    The Big Bands and Swing (1925--1945)}
    \vspace{1cm}
\end{center}

\tableofcontents
\newpage

% ============================================================
\section{Knowledge Tracker}
% ============================================================

\textit{Move items between sections as you study. The goal is to get everything into Green.}

\subsection*{\textcolor{greenzone}{Green --- Know From Memory}}
\begin{itemize}[label=$\square$]
    \item (move items here as you master them)
\end{itemize}

\subsection*{\textcolor{yellowzone}{Yellow --- Partial / Need Notes}}
\begin{itemize}[label=$\square$]
    \item Fletcher Henderson = innovated 15-piece big band format ($\sim$1924): 3 trumpets, 3 trombones, 5 saxes, 4 rhythm (piano, bass, drums, guitar)
    \item Duke Ellington coined ``swing'' via ``It Don't Mean a Thing If It Ain't Got That Swing'' (1931 composition, 1932 recording)
    \item ``Ellington effect'' = coined by Billy Strayhorn (1940); moaning, sustaining passages underneath trumpet melodies
    \item Contrafact = writing a new melody over a pre-existing chord progression to create a new composition
    \item ``Rhythm changes'' = chord progression from Gershwin's ``I Got Rhythm''; most commonly used progression in jazz alongside the blues
\end{itemize}

\subsection*{\textcolor{redzone}{Red --- Didn't Know / Got Wrong}}
\begin{itemize}[label=$\square$]
    \item Cotton Club: Ellington was resident composer in Harlem in the 1930s; all performers were Black, all patrons were white
    \item ``Jungle style'' = pejorative 1930s term for Ellington's West African-inspired sound (plunger mutes, growling brass)
    \item Papa Joe Jones moved four-on-the-floor from bass drum to hi-hat, creating lighter swing feel; influenced bebop
    \item Benny Goodman = ``King of Swing'' (controversial; Black bandleaders overlooked); one of first to integrate his band (hired Lionel Hampton, Teddy Wilson, Charlie Christian, Cootie Williams)
    \item Saxophone replaced clarinet as most prominent wind instrument in jazz during swing era
\end{itemize}

\subsection*{\textcolor{grayzone}{Not Yet Tested}}
\begin{itemize}[label=$\square$]
    \item Duke Ellington (1899--1974): born Washington, DC, upper-middle-class, classically trained; led orchestra 50 years (1923--1974)
    \item Count Basie (1904--1984): born New Jersey; took over Benny Moten's band in Kansas City after Moten died (1935)
    \item Head arrangements = compositions created spontaneously in rehearsal ``from the top of your head''
    \item Riff = short melodic phrase that's repeated; Basie placed much greater emphasis on riffs
    \item Basie = minimalist (opposite of Ellington's elaborate compositions and harmonies)
    \item All-American Rhythm Section (1937--1947): Basie (piano), Papa Joe Jones (drums), Walter Page (bass), Freddie Green (guitar)
    \item Four on the floor = bass drum hitting every quarter note; originated as ``Big Four'' in New Orleans jazz
    \item Billy Strayhorn wrote ``Take the A Train'' (not Ellington); became Ellington's theme song in the 1940s
    \item Ellington's five compositional functions: dance songs, jungle style (Cotton Club), mood pieces, pop tunes, abstract compositions
    \item Big Three tenor saxophonists: Coleman Hawkins (Henderson's band), Lester Young (Basie's band), Ben Webster (Ellington's band)
    \item Johnny Hodges = Ellington's lead alto saxophonist; dominant voice on alto until Charlie Parker in 1940s
    \item 32-bar AABA form = standard popular song form (B section = bridge)
    \item Billie Holiday (1915--1959): voice like an instrument/horn; exchanged nicknames with Lester Young (``Lady Day'' and ``Prez'')
    \item Ella Fitzgerald (1917--1996): ``First Lady of Song''; four-octave range; master of scat singing
    \item Coleman Hawkins = ``Father of the Jazz Tenor Saxophone''; vertical approach (playing up and down chord changes)
    \item Lester Young = ``cool'' sound; horizontal/melodic approach (contrast with Hawkins' ``hot''/vertical style)
    \item Swing music spread via radio, movies, and recordings; jazz WAS popular music in the 1930s because it was dance music
    \item Lindy Hop = popular swing-era dance style
    \item Harlem Renaissance saw Ellington as ``New School'' of Black Americans vs.\ Armstrong as ``old school''
    \item LeRoi Jones (Amiri Baraka), \textit{Blues People} (1963) = one of first books by an African American scholar on jazz history
    \item Gunther Schuller wrote \textit{The Swing Era} (companion to \textit{Early Jazz})
    \item ``Strange Fruit'' (1939) = one of first protest songs in jazz, about lynching; written by Abel Meeropol; Holiday's record company tried to ban it
\end{itemize}

\newpage

% ============================================================
\section{The Era: Big Bands and Swing (1925--1945)}
% ============================================================

\subsection{What This Movement Contributes}
\begin{enumerate}
    \item \textbf{Big Band Format} --- Fletcher Henderson created the 15-piece template ($\sim$1924) that remains standard today
    \item \textbf{Swing as Dance Music} --- jazz became America's popular music; people danced to it (Lindy Hop)
    \item \textbf{The Saxophone's Rise} --- saxophone overtook clarinet as the dominant jazz wind instrument
    \item \textbf{Compositional Sophistication} --- Ellington elevated jazz composition (mood pieces, abstract works, descriptive pieces)
    \item \textbf{The Soloist Within the Ensemble} --- star soloists emerged from big bands (Hawkins, Young, Webster, Hodges)
    \item \textbf{Contrafact Tradition} --- jazz musicians began building new compositions on existing chord progressions (``rhythm changes'')
\end{enumerate}

\subsection{Era Characteristics}
\begin{itemize}
    \item \textbf{Rhythmic Feel:} Swing rhythm; four-on-the-floor (bass drum or hi-hat); beats 2 and 4
    \item \textbf{Ensemble Style:} Large bands (15 pieces); section-based writing (brass vs.\ reeds); call-and-response between sections
    \item \textbf{Improvisation:} Extended solos within arranged frameworks; head arrangements (Basie) vs.\ elaborate compositions (Ellington)
    \item \textbf{Key Forms:} 32-bar AABA (popular songs, ``rhythm changes''); 12-bar blues still prominent
    \item \textbf{Technology:} Radio broadcasts, sound film (early 1930s), and recordings spread swing nationally and globally
    \item \textbf{Sound:} Section playing (saxes in unison, brass in unison); ``hot'' vs.\ ``cool'' soloist approaches
\end{itemize}

\subsection{Key Concepts}
\begin{itemize}
    \item \textbf{Big Band Format:} 3 trumpets, 3 trombones, 5 saxophones, 4 rhythm (piano, bass, drums, guitar)
    \item \textbf{Ellington Effect:} Moaning, sustaining accompaniment passages underneath trumpet melodies (coined by Billy Strayhorn)
    \item \textbf{Jungle Style:} Pejorative 1930s term for Ellington's West African-inspired sound; plunger mutes, growling brass (Bubber Miley)
    \item \textbf{Head Arrangement:} Composition created spontaneously during rehearsal; melody often emerges last
    \item \textbf{Riff:} Short melodic phrase that repeats; central to Basie's approach
    \item \textbf{Contrafact:} New melody written over pre-existing chord progression (e.g., ``Cottontail'' uses ``I Got Rhythm'' changes)
    \item \textbf{Rhythm Changes:} Chord progression from Gershwin's ``I Got Rhythm''; most commonly used progression in jazz alongside blues
    \item \textbf{Four on the Floor:} Bass drum (or hi-hat) on every quarter note; Papa Joe Jones moved it from bass drum to hi-hat
    \item \textbf{Hot vs.\ Cool:} ``Hot'' (Hawkins, Webster) = aggressive, vertical, influenced by Armstrong; ``Cool'' (Young) = relaxed, horizontal, melodic
    \item \textbf{Vertical vs.\ Horizontal Improvisation:} Vertical = playing up and down chord tones (Hawkins); Horizontal = playing melodic lines across chords (Young)
\end{itemize}

\newpage

% ============================================================
\section{Artists and Songs --- Syllabus Listening}
% ============================================================

% --- TAKE THE A TRAIN ---
\subsection{Duke Ellington --- ``Take the A Train'' (1941)}

\textbf{Era:} Swing / Big Band (1940s)

\subsubsection*{Key Facts}
\begin{enumerate}
    \item Written by Billy Strayhorn (NOT Ellington); became Ellington's theme song in the 1940s
    \item Strayhorn was Ellington's pianist and composer-collaborator
    \item Features the ``Ellington effect'' --- moaning, sustaining passages underneath the melody
    \item Notable use of dynamics throughout the arrangement
    \item Demonstrates Ellington's ability to showcase band as a cohesive orchestral unit
    \item 32-bar AABA form
\end{enumerate}

\textbf{Instrumentation:} Full big band --- trumpet section, trombone section, saxophone section, rhythm section (piano, bass, drums, guitar)\\
\textbf{Song Form:} 32-bar AABA

\listenbox{Catchy opening melody (trumpet), Ellington effect (sustained sax/trombone passages under melody), dynamic contrasts (loud/soft), full band vs.\ solo sections, polished arrangement, swinging rhythm section}

\bigskip
\textcolor{rulecolor}{\rule{\textwidth}{0.2pt}}

% --- COTTONTAIL ---
\subsection{Duke Ellington --- ``Cottontail'' (1940)}

\textbf{Era:} Swing / Big Band (1940s)

\subsubsection*{Key Facts}
\begin{enumerate}
    \item Based on ``I Got Rhythm'' chord changes --- a \textbf{contrafact}
    \item Features tenor saxophonist Ben Webster as primary soloist
    \item Ben Webster exemplifies the ``hot sound'' on saxophone, influenced by Louis Armstrong's hot trumpet style
    \item Webster was one of the Big Three tenor saxophonists (with Hawkins and Young)
    \item Demonstrates how jazz musicians built new compositions on pre-existing chord progressions
    \item 32-bar AABA form (same structure as ``I Got Rhythm'')
\end{enumerate}

\textbf{Instrumentation:} Full big band with Ben Webster (tenor sax) as featured soloist\\
\textbf{Song Form:} 32-bar AABA (rhythm changes / ``I Got Rhythm'' chord progression)

\listenbox{Ben Webster's ``hot'' tenor sax solo (aggressive, breathy, Armstrong-influenced), new melody over familiar chord changes, call-and-response between sax section and brass section, energetic swing feel, arranged ensemble passages between solos}

\bigskip
\textcolor{rulecolor}{\rule{\textwidth}{0.2pt}}

% --- ONE O'CLOCK JUMP ---
\subsection{Count Basie --- ``One O'Clock Jump'' (1937)}

\textbf{Era:} Swing / Big Band / Kansas City Jazz (late 1930s)

\subsubsection*{Key Facts}
\begin{enumerate}
    \item A blues-based \textbf{head arrangement} --- created spontaneously during rehearsal, not from a written score
    \item The melody doesn't appear until near the very end of the piece (characteristic of head arrangements)
    \item Features Lester Young on tenor saxophone with a ``cool'' sound (contrast with Ben Webster's ``hot'' sound)
    \item Demonstrates Basie's minimalist piano style
    \item Papa Joe Jones moves four-on-the-floor from bass drum to hi-hat, creating lighter, more buoyant feeling
    \item Freddie Green's guitar provides the underlying four-feel rhythm
    \item Texture variety: band drops out during quieter solos, Papa Joe Jones switches between ride cymbal and hi-hat
\end{enumerate}

\textbf{Instrumentation:} Full big band --- All-American Rhythm Section: Basie (piano), Papa Joe Jones (drums), Walter Page (bass), Freddie Green (guitar); plus horn sections with Lester Young (tenor sax) featured\\
\textbf{Song Form:} 12-bar blues; head arrangement structure (solos first, melody last)

\listenbox{Lester Young's ``cool'' relaxed tenor sax tone, Basie's spare/minimalist piano comping, hi-hat timekeeping (not bass drum), Freddie Green's steady guitar pulse, texture changes (band drops out for piano solo, swells for ensemble), riff-based melody arriving late in the arrangement, blues form throughout}

\bigskip
\textcolor{rulecolor}{\rule{\textwidth}{0.2pt}}

% --- OH LADY BE GOOD ---
\subsection{Count Basie / Lester Young --- ``Oh! Lady Be Good'' (1936)}

\textbf{Era:} Swing / Big Band / Kansas City Jazz (mid-1930s)

\subsubsection*{Key Facts}
\begin{enumerate}
    \item Features Lester Young as primary soloist --- quintessential example of his ``cool'' approach
    \item Lester Young's horizontal/melodic improvisation style contrasts with Coleman Hawkins' vertical/chordal approach
    \item Young played in Count Basie's band; one of the Big Three tenor saxophonists
    \item Young's relaxed, laid-back phrasing influenced the ``cool jazz'' movement of the 1950s
    \item Exchanged nicknames with Billie Holiday: she called him ``Prez'' (president of the saxophone), he called her ``Lady Day''
\end{enumerate}

\textbf{Instrumentation:} Small group or big band with Lester Young (tenor sax) featured\\
\textbf{Song Form:} 32-bar AABA (standard popular song form)

\listenbox{Lester Young's light, airy, ``cool'' tone (contrast with Hawkins' heavier vibrato), melodic/horizontal approach to improvisation (lines flow across chord changes rather than outlining individual chords), relaxed behind-the-beat phrasing, smooth legato lines, minimal vibrato compared to Hawkins}

\bigskip
\textcolor{rulecolor}{\rule{\textwidth}{0.2pt}}

% --- BODY AND SOUL ---
\subsection{Coleman Hawkins --- ``Body and Soul'' (1939)}

\textbf{Era:} Swing (late 1930s)

\subsubsection*{Key Facts}
\begin{enumerate}
    \item Considered the first great jazz tenor saxophone solo
    \item Hawkins = ``Father of the Jazz Tenor Saxophone''; nicknames ``Hawk'' and ``Bean'' (for being intellectual)
    \item Features mostly improvisation with only brief quotation of the original melody
    \item Demonstrates \textbf{vertical approach} to improvisation --- playing up and down chord changes
    \item Hawkins emerged from Fletcher Henderson's band (1923--1934), spent 1934--1939 in London, recorded this upon return
    \item Extended saxophone technique into the altissimo register (extreme high range)
    \item Extensive vibrato throughout notes; smooth tone with jumping quality and texture
    \item Tenor saxophone often considered the instrument closest to the human voice
    \item Hawkins was pivotal in the transition from swing to bebop; reportedly the first to hire Thelonious Monk
\end{enumerate}

\textbf{Instrumentation:} Tenor saxophone (Hawkins) with rhythm section (piano, bass, drums)\\
\textbf{Song Form:} 32-bar AABA; Hawkins largely abandons the melody and improvises over the chord changes

\listenbox{Almost entirely improvised (melody barely stated), vertical/chordal improvisation (arpeggios moving through chord tones), extensive vibrato, altissimo register passages, smooth but active melodic movement, romantic ballad interpretation, lush tone quality}

\bigskip
\textcolor{rulecolor}{\rule{\textwidth}{0.2pt}}

% --- SING SING SING ---
\subsection{Benny Goodman --- ``Sing, Sing, Sing'' (1937)}

\textbf{Era:} Swing (late 1930s)

\subsubsection*{Key Facts}
\begin{enumerate}
    \item One of the most iconic pieces in jazz history
    \item Goodman (1909--1986) = ``King of Swing''; born in Chicago to Jewish immigrant parents from Poland
    \item Classically trained clarinetist who started playing professionally at age 14
    \item Unique as a drum feature centered around Gene Krupa's drumming
    \item Over seven minutes long, split across two sides of a record (unusual for the three-minute standard)
    \item The drum rhythm is one of the most frequently quoted elements in American music
    \item Fletcher Henderson was Goodman's arranger (same Henderson who worked with Armstrong and Hawkins)
    \item ``King of Swing'' title was controversial --- Black bandleaders (Ellington, Basie, Henderson, Cab Calloway) were overlooked
    \item Goodman was one of the first white bandleaders to integrate his band: hired Teddy Wilson (piano), Lionel Hampton (vibraphone), Charlie Christian (guitar), Cootie Williams (trumpet)
    \item Goodman's clarinet: beautiful classical tone combined with blues feeling and active melodic movement
\end{enumerate}

\textbf{Instrumentation:} Full big band with Gene Krupa (drums) and Benny Goodman (clarinet) as featured soloists\\
\textbf{Song Form:} Extended form (unusual for the era); built around recurring drum patterns

\listenbox{Gene Krupa's driving tom-tom rhythm (the iconic drum pattern), extensive call-and-response between drums and horns, Goodman's clarinet solo (classical tone + blues feeling), extended length (7+ minutes), building energy and intensity, drum-dominated texture (unusual for jazz)}

\bigskip
\textcolor{rulecolor}{\rule{\textwidth}{0.2pt}}

% --- GOD BLESS THE CHILD ---
\subsection{Billie Holiday --- ``God Bless the Child'' (1942)}

\textbf{Era:} Swing (early 1940s)

\subsubsection*{Key Facts}
\begin{enumerate}
    \item Holiday's most famous self-composed song
    \item Billie Holiday (1915--1959): died young at age 44
    \item Voice was like an instrument/horn --- fragile and distinct timbre, relatively small vocal range but immense feeling
    \item Vocal characteristics: very legato phrasing, words spilling into each other, unique timbre, yearning quality, wide vowel placement
    \item Exchanged nicknames with Lester Young: she was ``Lady Day,'' he was ``Prez''
    \item Catapulted jazz vocals to new level of fame
    \item Raw emotional quality and unique tone distinguished her from other vocalists
    \item Use of dynamics: quieter moments building to louder endings
    \item Voice described as woeful, wistful, reminiscent, sober
    \item 32-bar ABA form
\end{enumerate}

\textbf{Instrumentation:} Voice (Holiday) with small ensemble accompaniment\\
\textbf{Song Form:} 32-bar ABA

\listenbox{Legato phrasing (words flow into each other), unique/fragile vocal timbre (horn-like quality), yearning emotional quality, dynamic contrast (quiet to loud), behind-the-beat phrasing, wide vowel placement, minimal but effective vibrato, storytelling delivery}

\bigskip
\textcolor{rulecolor}{\rule{\textwidth}{0.2pt}}

% --- BLUE SKIES ---
\subsection{Ella Fitzgerald --- ``Blue Skies'' (1958)}

\textbf{Era:} Swing / Post-Swing (late 1950s)\\
\textit{from Ella Fitzgerald Sings the Irving Berlin Songbook}

\subsubsection*{Key Facts}
\begin{enumerate}
    \item Ella Fitzgerald (1917--1996) = ``First Lady of Song''
    \item Four-octave vocal range; very clear, big voice
    \item Considered the master of scat singing (though Louis Armstrong created the technique)
    \item Uses chromaticism and runs up and down chords --- vertical approach similar to Coleman Hawkins
    \item Incorporates humor through quotation in her scatting (quotes the bridal march, Gershwin tunes, etc.)
    \item Approach contrasts with Holiday's: Fitzgerald = virtuosic/vertical (like Hawkins); Holiday = emotional/intimate (like Webster in ballads)
\end{enumerate}

\textbf{Instrumentation:} Voice (Fitzgerald) with ensemble/big band accompaniment\\
\textbf{Song Form:} 32-bar AABA (standard Irving Berlin song form); includes extended scat sections

\listenbox{Scat singing (improvised syllables), chromaticism and rapid melodic runs, four-octave range on display, quotation of other melodies within scat solos (humor), clear and powerful tone, rhythmic precision, contrast with Holiday's fragile intimacy --- Fitzgerald is virtuosic and exuberant}

\newpage

% ============================================================
\section{Additional Pieces Discussed in Class}
% ============================================================

\textit{Not on the syllabus listening list but discussed in lecture --- may appear on quizzes.}

\bigskip

% --- IT DON'T MEAN A THING ---
\subsection{Duke Ellington --- ``It Don't Mean a Thing If It Ain't Got That Swing'' (1932)}

\begin{enumerate}
    \item Gave the entire swing movement its name
    \item Composed 1931, recorded 1932
    \item 32-bar AABA song form
    \item Features elaborate singing and an alto saxophone solo by Johnny Hodges
    \item Hodges was Ellington's lead alto saxophonist and dominant voice on alto until Charlie Parker
\end{enumerate}

\listenbox{AABA form, Johnny Hodges' alto sax solo, vocal scatting, the word ``swing'' used to define the rhythmic feel of an entire era}

\bigskip
\textcolor{rulecolor}{\rule{\textwidth}{0.2pt}}

% --- DAYBREAK EXPRESS ---
\subsection{Duke Ellington --- ``Daybreak Express'' (1934)}

\begin{enumerate}
    \item Abstract composition deliberately mimicking the sound of a train on tracks
    \item Uses chromatic scales to create intensity and the Doppler effect
    \item Demonstrates Ellington's compositional ambition beyond dance music --- a descriptive/programmatic piece
\end{enumerate}

\listenbox{Train sounds/imagery, chromatic scales, building intensity, Doppler-like effects, orchestral approach to jazz composition}

\bigskip
\textcolor{rulecolor}{\rule{\textwidth}{0.2pt}}

% --- EAST ST LOUIS TOODLE-OO ---
\subsection{Duke Ellington --- ``East St. Louis Toodle-oo'' (1927)}

\begin{enumerate}
    \item Best example of the ``Ellington effect''
    \item Features Bubber Miley's trumpet with moaning saxophone and tuba accompaniment underneath
    \item Miley used a plunger mute and growling technique = ``jungle style''
\end{enumerate}

\listenbox{Plunger-muted trumpet (Bubber Miley), growling brass, moaning sustained accompaniment under trumpet melody, dark timbral quality}

\bigskip
\textcolor{rulecolor}{\rule{\textwidth}{0.2pt}}

% --- STRANGE FRUIT ---
\subsection{Billie Holiday --- ``Strange Fruit'' (1939)}

\begin{enumerate}
    \item One of the first protest songs in jazz history --- about lynching in America
    \item Written by Abel Meeropol (not Holiday)
    \item Record company tried to ban her from performing it; she refused and continued, leading to jail time
    \item Sung in minor key, straightforward delivery, no vibrato, confrontational
    \item Ominous piano introduction; meant to make audience uncomfortable
    \item Matter-of-fact tone was mocking and accusatory
\end{enumerate}

\listenbox{Minor key, ominous piano intro, absence of vibrato (unlike her other singing), stark/confrontational delivery, uncomfortable silences, dramatic contrast with typical swing-era entertainment music}

\bigskip
\textcolor{rulecolor}{\rule{\textwidth}{0.2pt}}

% --- RUBY MY DEAR ---
\subsection{Coleman Hawkins --- ``Ruby My Dear'' (1957, with Thelonious Monk)}

\begin{enumerate}
    \item Showcases Hawkins' romantic ballad style later in his career
    \item Stays within chord tones; demonstrates unbelievable tone quality
    \item Hawkins was supportive of new music and reportedly the first to hire Monk
\end{enumerate}

\listenbox{Romantic ballad interpretation, rich warm tone, vibrato, chord-tone-based improvisation, interplay with Monk's angular piano}

\bigskip
\textcolor{rulecolor}{\rule{\textwidth}{0.2pt}}

% --- TIME ON MY HANDS ---
\subsection{Ben Webster --- ``Time On My Hands'' (mid-1950s)}

\begin{enumerate}
    \item Webster known more as a balladeer than uptempo player
    \item Vocal quality with crying sound and extensive use of air/breath in tone
    \item Very early and fast vibrato adding yearning quality
    \item Album \textit{Soulville} recommended for more of this style
\end{enumerate}

\listenbox{Breathy/airy tone, crying quality, fast vibrato, extreme tenderness and nuance, vocal-like phrasing, dynamic subtlety}

\newpage

% ============================================================
\section{Listening Response Template}
% ============================================================

Name, composer, and year are PROVIDED on the quiz. Focus on what you HEAR. Aim for 5 observations:

\begin{enumerate}
    \item \textbf{Instrumentation/Ensemble} --- Big band or small group? Which sections (brass, reeds, rhythm)? Who is the soloist?
    \item \textbf{Era/Style} --- Swing, Kansas City, big band? What clues? (full band = swing era; riff-based = Basie; lush orchestration = Ellington; drum feature = Goodman)
    \item \textbf{Rhythm/Feel} --- Swinging? Four-on-the-floor (bass drum or hi-hat)? Riff-based? Light/heavy?
    \item \textbf{Melody/Improvisation} --- Head arrangement or composed? Contrafact? Vertical or horizontal soloist? Scat singing? Hot or cool?
    \item \textbf{Texture/Form/Special} --- 32-bar AABA or 12-bar blues? Call-and-response between sections? Ellington effect? Dynamic contrasts? Plunger mute? Extended length?
    \item \textbf{Emotion/Context} --- Dance music or art music? Protest? How does it connect to readings?
\end{enumerate}

\noindent\fcolorbox{rulecolor}{lightbg}{\parbox{\dimexpr\textwidth-2\fboxsep-2\fboxrule}{\small\textbf{Template:} ``This piece features [ensemble type] with [soloist] as the primary voice. The style is [era/style] based on [clues]. The form is [form], evidenced by [what you hear]. A notable feature is [detail], which connects to [course concept].''}}

% ============================================================
\section{General Understanding}
% ============================================================

\subsection{Movement's Place in History}
\begin{itemize}
    \item Big band format (Henderson, $\sim$1924) standardized jazz into large ensemble music
    \item Swing era = jazz as America's popular music; people danced to it
    \item Technology (radio, film, recordings) spread swing nationally and globally
    \item Saxophone overtook clarinet as dominant wind instrument; planted seeds for modern jazz
    \item Star soloists emerged from big bands, setting the stage for small-group bebop in the 1940s
\end{itemize}

\subsection{Cultural Context}
\begin{itemize}
    \item Cotton Club: Ellington's creative incubator, but with racist dynamic (Black performers, white-only audience)
    \item ``King of Swing'' controversy: Benny Goodman (white) received title while Black innovators were overlooked
    \item Goodman also one of first to integrate his band --- hired Hampton, Wilson, Christian, Williams despite criticism
    \item Harlem Renaissance: Ellington = ``New School'' of Black Americans; Armstrong = ``old school''
    \item 1930s paradox: African American music celebrated while African Americans themselves were not embraced by society
    \item LeRoi Jones (Amiri Baraka), \textit{Blues People} (1963): one of first African American scholarly treatments of jazz history
\end{itemize}

\subsection{Musical Evolution}
\begin{itemize}
    \item Ellington: elaborate compositions, orchestral ambition, five compositional functions, ``Ellington effect''
    \item Basie: minimalist, head arrangements, riff-based, lighter swing (hi-hat instead of bass drum)
    \item Contrafact tradition: new melodies over existing chord progressions became fundamental jazz practice
    \item Hot vs.\ Cool dichotomy: Hawkins/Webster (hot, vertical) vs.\ Young (cool, horizontal) --- this tension carries into bebop and cool jazz
    \item Papa Joe Jones' hi-hat innovation directly influenced bebop drumming in the 1940s
\end{itemize}

\subsection{Key Tensions}
\begin{itemize}
    \item \textbf{Composed vs.\ Spontaneous:} Ellington's elaborate scores vs.\ Basie's head arrangements
    \item \textbf{Hot vs.\ Cool:} Hawkins/Webster's aggressive approach vs.\ Young's laid-back melodicism
    \item \textbf{Entertainment vs.\ Art:} Swing was dance music, but Ellington pushed toward abstract composition (``Daybreak Express'')
    \item \textbf{Black Innovation, White Recognition:} Henderson arranged for Goodman; Goodman got the ``King of Swing'' title
    \item \textbf{Virtuosity vs.\ Emotion:} Fitzgerald's four-octave scat vs.\ Holiday's fragile intimacy
\end{itemize}

\end{document}
